\documentclass[11pt,a4paper]{article}
\usepackage[utf8]{inputenc}
\usepackage[T1]{fontenc}
\usepackage[french]{babel}
\usepackage{amsmath}
\usepackage{amssymb}
\usepackage{listings}
\usepackage{xcolor}
\usepackage{graphicx}

% Configuration pour les blocs de code VHDL
\lstset{
    language=VHDL,
    basicstyle=\ttfamily\footnotesize,
    keywordstyle=\color{blue}\bfseries,
    commentstyle=\color{green!60!black},
    stringstyle=\color{red},
    showstringspaces=false,
    frame=single,
    numbers=left,
    numberstyle=\tiny\color{gray},
    breaklines=true
}

\title{TD 1 : Exercices simples de synthèse logique \\ \large EE341 : Architecture matérielle avancée 1}
\author{ESISAR-INPG-E13}
\date{}

\begin{document}

\maketitle

\section*{Exercice 1.1}
On considère le programme ci-dessous : [cite: 284, 285]
\begin{enumerate}
    \item Déduire de ce programme, par une construction méthodique, un schéma (bascules et portes logiques). [cite: 286, 287]
    \item Compléter le chronogramme Fig 1.1. [cite: 288]
\end{enumerate}

\begin{lstlisting}
entity transitm is
port (hor, e : in bit;
      s      : out bit);
end transitm;
architecture quasi_struct of transitm is
signal qa, qb : bit;
begin
    s <= qa xor qb;
    schem: process (hor)
    begin
        if hor'event and hor = '1' then
            qa <= e;
            qb <= qa;
        end if;
    end process schem;
end quasi_struct;
\end{lstlisting}

\subsection*{Correction}
\textbf{1. Schéma :}\\
Le schéma correspondant se trouve dans le fichier SVG 1.1 (à inclure ici). Il est constitué de deux bascules D en cascade pour \texttt{qa} et \texttt{qb}, pilotées par la même horloge \texttt{hor}, dont les sorties attaquent une porte XOR pour générer \texttt{s}. [cite: 312, 313, 314, 315, 316, 317, 318, 319, 320, 321, 322, 323]

\section*{Exercice 1.2}
On considère le programme ci-dessous : [cite: 327, 328]
\begin{enumerate}
    \item Compléter le code manquant. [cite: 329]
    \item Déduire un schéma (bascules et portes logiques) de ce programme. [cite: 330]
    \item Caractériser les sorties : synchrones ou asynchrones ? [cite: 331]
    \item Compléter le chronogramme Fig 1.2. [cite: 332]
\end{enumerate}

\begin{lstlisting}
entity condit is
port (
    clk    : in bit;
    eclk   : in bit;
    e      : in bit;
    s_clk  : out bit;
    s_eclk : out bit);
end condit;
architecture reg_lvl of condit is
    signal qe : bit ;
begin
    s_clk <= e and not (qe);
    
    echl : process (clk)
    begin
        if clk'event and clk = '1' then
            qe <= e;
        end if;
    end process echl;
    
    ech : process (eclk)
    begin
        if eclk'event and eclk = '1' then
            s_eclk <= (qe and not(e));
        end if;
    end process ech;
end reg_lvl;
\end{lstlisting}

\subsection*{Correction}
\textbf{2. Schéma :}\\
Voir le fichier SVG 1.2 pour le schéma de synthèse. [cite: 373, 374, 375, 376, 377, 378, 379, 380, 381, 382]

\textbf{3. Caractérisation des sorties :}\\
La sortie \texttt{s\_clk} est \textbf{asynchrone} car elle est affectée en dehors de tout processus synchrone. [cite: 351, 385, 386]\\
La sortie \texttt{s\_eclk} est \textbf{synchrone} (synchronisée sur \texttt{eclk}) car son affectation se trouve dans un processus sensible au front montant de l'horloge \texttt{eclk}. [cite: 361, 362, 387, 388]

\section*{Exercice 1.3}
On considère le programme ci-dessous : [cite: 394, 395]
\begin{enumerate}
    \item Ce code est-il correct ? [cite: 396]
    \item Déduire un schéma (bascules et portes logiques) de ce programme. [cite: 397]
    \item Compléter le chronogramme Fig 1.3. [cite: 397]
\end{enumerate}

\begin{lstlisting}
entity mise_feu is
port (
    clk : in bit;
    ina : in bit;
    maf : out bit);
end mise_feu;
architecture detect of mise_feu is
    signal maf_sync : bit;
begin
    maf <= maf_sync;
    sync : process (clk)
    begin
        if clk'event and clk = '1' then
            maf_sync <= ina and not maf_sync;
        end if;
    end process sync;
end detect;
\end{lstlisting}

\subsection*{Correction}
\textbf{2. Schéma :}\\
Voir le fichier SVG 1.3 (boucle de rétroaction sur la bascule avec porte AND). [cite: 405, 406, 407, 408, 409, 410]

\section*{Exercice 2.1}
On considère la fonction logique de la figure 2.1 (Bloc fonctionnel) : [cite: 432, 433, 448]
\begin{enumerate}
    \item Compléter le programme pour réaliser cette fonction. [cite: 434]
    \item Réaliser un chronogramme de test de cette fonction (compléter le chronogramme fig 2.2). [cite: 436]
\end{enumerate}

\subsection*{Correction}
\textbf{1. Code VHDL :} [cite: 449-467]
\begin{lstlisting}
entity detect is
port (
    clk : in bit;
    e   : in bit;
    s   : out bit);
end detect;
architecture reg of detect is
    signal qe : bit;
begin
    sync : process (clk)
    begin
        if clk'event and clk = '1' then
            qe <= e;
            s  <= qe xor e;
        end if;
    end process sync;
end reg;
\end{lstlisting}
*(Le schéma est illustré dans le SVG 2.1)* [cite: 437, 438, 439, 440, 441, 442, 443, 444, 445, 446, 447]

\section*{Exercice 2.2}
On considère l'équation logique ci-dessous : [cite: 480, 483]
$$Q_{(t+1)} = \overline{(init \cdot (T \oplus Q_{(t)}))}$$ [cite: 484]
\begin{enumerate}
    \item Concevoir le schéma de réalisation de cette équation. [cite: 485]
    \item Compléter le code VHDL pour réaliser cette fonction. [cite: 486]
    \item Le cahier des charges évolue : il faut transformer le signal \texttt{init} en une commande << actif bas >> (si $init='0'$), de remise à zéro asynchrone. [cite: 487, 488]
    \item Transformer le programme pour qu'il utilise le type \texttt{std\_logic}. Rajouter une commande de haute impédance << \texttt{oe} >> active basse, qui pilote la sortie uniquement lorsque << \texttt{oe} >> est actif. [cite: 489, 490]
\end{enumerate}

\subsection*{Correction}

\textbf{1. Schéma :}\\
Voir le fichier SVG 2.2. Il est composé d'une porte XOR, d'une porte NAND et d'une bascule D. [cite: 499, 500, 501, 502, 503, 504, 505, 506, 507, 508]

\textbf{2. Code VHDL classique :} [cite: 509-524]
\begin{lstlisting}
entity basc is
port (T, clk, init : in bit;
      s            : out bit);
end basc;
architecture primitive of basc is
    signal q : bit;
begin
    s <= q;
    process (clk)
    begin
        if clk'event and clk = '1' then
            q <= not (init and (T xor q));
        end if;
    end process;
end primitive;
\end{lstlisting}

\textbf{3. Programme avec la commande init en remise à zéro asynchrone :} [cite: 530-549]
\begin{lstlisting}
entity basc is
port (T, clk, init : in bit;
      s            : out bit);
end basc;
architecture primitive of basc is
    signal q : bit;
begin
    s <= q;
    process (clk, init)
    begin
        if init = '0' then
            q <= '0';
        elsif clk'event and clk = '1' then
            q <= not (T xor q);
        end if;
    end process;
end primitive;
\end{lstlisting}

\textbf{4. Programme avec le type \texttt{std\_logic} et \texttt{oe} :} [cite: 550-569]
\begin{lstlisting}
library IEEE;
use IEEE.STD_LOGIC_1164.ALL;

entity basc is
port (
    T, clk, init, oe : in std_logic;
    s                : out std_logic);
end basc;
architecture primitive of basc is
    signal q : std_logic;
begin
    s <= q when oe = '0' else 'Z';
    process (clk, init)
    begin
        if init = '0' then
            q <= '0';
        elsif clk'event and clk = '1' then
            q <= not (T xor q);
        end if;
    end process;
end primitive;
\end{lstlisting}

\end{document}